% Executive Summary — Chapter-level
%
% Written by the lead coauthor after individual contributions are drafted.
% Synthesizes top-line findings across all coauthors. Keep to 1–2 pages.

\section*{Executive Summary}
\addcontentsline{toc}{section}{Executive Summary}

[Chapter-level synthesis. This section is typically the last thing written —
it pulls together the Key Messages from each coauthor's contribution into a
coherent chapter-level narrative. It should:]

\begin{itemize}
  \item State the chapter's central question or scope in one paragraph.
  \item Highlight the most important findings across coauthor contributions, with
        confidence labels (e.g., ``Climate-driven heat exposure reduces labor
        supply in outdoor industries; \confidence{high}'').
  \item Identify where the chapter's findings update, extend, or contest prior
        assessments (NCA5, IPCC AR6).
  \item Flag the major open questions the chapter could not resolve.
\end{itemize}

\subsection*{Cross-Cutting Key Messages}

[Optional: a numbered list of 3–5 chapter-level KMs that span coauthor
contributions. Individual KMs live inside each contribution; this section
captures themes that emerge across them.]

\begin{enumerate}
  \item \textbf{[Cross-cutting KM 1]} — [one-paragraph statement with
        confidence and likelihood labels.] \confidence{high} \likelihood{very likely}
  \item \textbf{[Cross-cutting KM 2]} — [...] \confidence{medium}
  \item \textbf{[Cross-cutting KM 3]} — [...] \confidence{high}
\end{enumerate}

\subsection*{Chapter Roadmap}

[One short paragraph: how the chapter is organized. Mention each coauthor
contribution by section title.]

\newpage
